\section{UART 历史、帧格式与 16550A}

\topic{初始化顺序}
固件先关闭 UART 中断，设置 DLAB，写入波特率除数 1，再关闭 DLAB 并配置
8 数据位、无校验、1 停止位。随后配置 FIFO 和 modem control。对 QEMU
模拟串口而言，宿主终端不真正按 115200 波特限制吞吐，但寄存器协议仍应按照
真实设备编写。

\begin{osgridlongtable}{p{0.24\textwidth}p{0.20\textwidth}p{0.46\textwidth}}
动作 & 端口/值 & 含义 \\ \hline
关闭中断 & 0x3F9 / 0x00 & 最早阶段采用轮询，不接收 IRQ \\ \hline
打开 DLAB & 0x3FB / 0x80 & 让偏移 0、1 指向除数锁存器 \\ \hline
设置除数 & 0x3F8 / 1，0x3F9 / 0 & 115200 基准波特率 \\ \hline
设置 8N1 & 0x3FB / 0x03 & 同时关闭 DLAB \\ \hline
设置 FIFO & 0x3FA / 0xC7 & 启用并清空 FIFO \\ \hline
设置 modem & 0x3FC / 0x0B & 拉起 DTR、RTS、OUT2 \\ \hline
\end{osgridlongtable}

\topic{发送不是直接写字节}
写数据前读取线路状态寄存器，等待 bit 5 表示发送保持寄存器为空。若直接连续
写而不检查，代码隐含了设备永远足够快的假设。当前实现用 16 位计数器进行最多
\code{0xFFFF} 次轮询，成功通过进位标志返回，超时清除进位标志。

\topic{两条日志的协议意义}
\code{RESET} 证明复位跳转、栈和 UART 初始化已经工作；
\code{SERIAL_READY} 证明完整字符串发送路径仍然可用。失败镜像在第一条消息
之后人为屏蔽就绪状态，用同一轮询函数走超时路径。因此测试既能识别“完全未
执行”，也能识别“执行了但设备没有完成”。

\topic{从电传打字机到 16550A}
串行通信把一个字符的各个位沿单条数据线依次发送，所需导线少于并行总线，
因而很早就用于连接终端、调制解调器和计算机。RS-232 定义了电气电平与连接
信号，UART 则在处理器可见的并行字节和线路上的串行比特之间转换。PC 保留的
COM 端口由这一历史接口演变而来。

早期 UART 每次只能暂存很少的数据，处理器服务不及时就会丢失字符。16550
系列加入 FIFO，使软件可以成批处理收发事件。现代机器即使不再暴露物理 DB-9
接口，固件和虚拟机仍广泛保留兼容 UART，因为它结构简单、初始化早、无需
图形栈，适合作为启动诊断通道。

\topic{异步帧在没有共享时钟时保持同步}
异步串口不在数据线上同时传送时钟。空闲线保持高电平，起始位拉低线路，接收
端由此确定采样相位，随后按约定波特率读取数据位、可选校验位和停止位。8N1
表示 8 个数据位、无校验、1 个停止位，一帧共 10 个比特时间。

\begin{pdfdiagram}[node distance=0mm]
  \node[os state,minimum width=14mm] (start) {起始\\0};
  \node[os block,right=of start,minimum width=46mm] (data) {数据位 0--7\\低位先行};
  \node[os hardware,right=of data,minimum width=14mm] (stop) {停止\\1};
  \draw[os arrow] (start) -- (data);
  \draw[os arrow] (data) -- (stop);
\end{pdfdiagram}
\diagramnote{115200 波特下，一个 8N1 字符占约 \(10/115200\) 秒，即 \(86.8\,\mu s\)。}

收发双方必须事先同意位宽、校验和速率。格式不一致时，字节值会被错误采样；
停止位缺失会在线路状态中形成帧错误。QEMU 把字符转交宿主终端时弱化了真实
时间约束，却保留了寄存器和状态机接口。

\topic{波特率除数的来源}
经典 PC UART 使用 1.8432 MHz 输入时钟，并在内部进行 16 倍采样。因此
波特率满足：
\[
  \mathrm{baud}=\frac{1\,843\,200}{16\times\mathrm{divisor}}.
\]
除数为 1 时得到 115200。除数是 16 位值，但 UART 只有 8 位数据端口，
所以低字节 DLL 与高字节 DLM 复用偏移 0 和 1。线路控制寄存器中的 DLAB
决定当前访问的是除数锁存器还是收发与中断寄存器。

这种寄存器复用来自早期有限端口资源和兼容性约束。它也意味着初始化不是任意
顺序的一组写操作：软件必须先置 DLAB，写完两个除数字节，再清 DLAB，随后
才能把偏移 0 当作发送保持寄存器。

\topic{初始化事务的状态轨迹}\par
\begin{osgridlongtable}{L{0.12\textwidth}L{0.18\textwidth}L{0.20\textwidth}L{0.34\textwidth}}
\textbf{序号} & \textbf{端口} & \textbf{写入值} & \textbf{新建立的设备状态}\\ \hline
\endfirsthead
\textbf{序号} & \textbf{端口} & \textbf{写入值} & \textbf{新建立的设备状态}\\ \hline
\endhead
1 & \code{0x3F9} & \code{0x00} & 关闭 UART 中断源\\ \hline
2 & \code{0x3FB} & \code{0x80} & DLAB 打开，偏移 0/1 指向除数\\ \hline
3 & \code{0x3F8} & \code{0x01} & 除数低字节为 1\\ \hline
4 & \code{0x3F9} & \code{0x00} & 除数高字节为 0\\ \hline
5 & \code{0x3FB} & \code{0x03} & DLAB 关闭，同时选择 8N1\\ \hline
6 & \code{0x3FA} & \code{0xC7} & 启用并清空 FIFO，设置触发级别\\ \hline
7 & \code{0x3FC} & \code{0x0B} & 设置 DTR、RTS 和 OUT2\\ \hline
\end{osgridlongtable}

表中的端口号和位模式来自硬件协议，所以源代码用具名常量表达，不能为了减少
行数直接散布数值。常量名同时记录项目、模块和功能；注释解释 DLAB 等状态
依赖，不重复“向端口写值”这种从指令本身即可看出的事实。

\topic{一个字节发送过程的双方状态机}
软件读取 LSR 的 THRE 位。位为 1 时，发送保持寄存器的所有权归软件，可以
写入下一个字节；写入后设备取得所有权并清除状态，待字节移入发送移位寄存器
后重新置位。轮询循环在这一所有权协议上等待。

当前过程把进位标志当作最小返回 ABI：置位表示字节已经交给 UART，清零表示
预算耗尽。调用者在发送字符串时逐字节检查该标志，所以任何一次超时都会终止
后续输出。后续进入 C++ 后可以用显式枚举表达成功、超时和线路错误，但汇编
与 C++ 边界仍需形成文档化 ABI。

\begin{failurebox}
THRE 只表示发送保持寄存器可接收新字节，不表示该字节已经完全离开串行线路。
需要确认全部发送完成时还要检查 TEMT。状态位名称相近，但对应不同硬件阶段。
\end{failurebox}

\topic{轮询与中断各自适用的阶段}
轮询实现短小，且在 IDT、PIC/APIC 和中断栈尚未建立时仍可使用，因此适合
复位后的第一条日志。缺点是等待期间独占处理器，吞吐和响应时间也依赖循环。
中断驱动在系统进入并发阶段后更高效，却要求中断路由、共享缓冲区和同步机制
已经可靠。

项目先使用有界轮询建立观测能力，后续再把串口迁移到中断驱动。这不是把临时
实现当成最终架构，而是让每一阶段只依赖已经验证的基础设施。
